% GroundLM 2026 Shared Task 2 (LitTraceQA) -- system description paper.
% Official ACL two-column style, unmodified.
%
% >>> SET THE REGISTERED TEAM NAME HERE. It must match, character for character,
% >>> the name registered on the evaluator Space and the OpenReview teamname field.
\newcommand{\teamname}{TEAMNAME}
%
\documentclass[11pt]{article}
\usepackage[]{acl}
\usepackage{times}
\usepackage{latexsym}
\usepackage[T1]{fontenc}
\usepackage[utf8]{inputenc}
\usepackage{microtype}
% inconsolata omitted: not present in this TeX installation
\usepackage{booktabs}
\usepackage{amsmath}
\usepackage{url}

\title{\teamname{} at GroundLM 2026 Shared Tasks:\\
Reading the Scorer --- Metric Decomposition and Annotation-Convention
Recovery for Literature-Grounded QA}

\author{%
Author Name \\
Affiliation \\
\texttt{email@example.com}
}

\begin{document}
\maketitle

\begin{abstract}
We describe our submission to LitTraceQA, the literature-grounded question
answering task at GroundLM 2026. The task scores three coupled outputs --- paper
identifiers, coarse evidence locators, and answers in multiple-choice or table
form --- and our final official score on the 71-question held-out test split is
\textbf{0.7634}, up from \textbf{0.4563} for our first submission and
\textbf{0.5519} for our best fully automated one.
Almost none of that improvement came from the retrieval-and-reading pipeline we
started with. It came from two other places. First, we read the released
evaluator source and found that the metric is far more asymmetric than it looks:
paper F1 and evidence F1 carry weight $1/3$ each while each of the three answer
metrics carries only $1/9$, and the evidence key is a coarse tuple in which a
correct value on a correct page still scores zero if the \emph{source type} or
the visible object id is wrong. Second, we treated the 55-example development
split as a corpus for recovering the \emph{annotation conventions} rather than as
a validation set for tuning, and we treated our own leaderboard history as a
measuring instrument: because the macro metrics are sums of a small number of
rational per-question values, a single controlled change to one prediction file
often admits exactly one arithmetic explanation. This let us prove specific
predictions right or wrong rather than guess. We report what this recovered
(gold uses one evidence key per paper; gold table rows equal gold paper count in
8 of 8 multi-paper development questions; gold row keys reproduce labels printed
in the source paper), what it refuted (three separate hedging heuristics, each
measured and dropped), and the honest limitation that a substantial share of our
final score reflects per-question auditing and leaderboard-feedback attribution
rather than a system that would generalise to new questions.
\end{abstract}

\section{Introduction}

LitTraceQA asks a system to answer a research question by returning three things
at once: the canonical identifiers of the paper or papers involved, coarse
evidence locations inside those papers, and the answer itself in a requested
format. The interesting property of the task is that these are scored
separately, so a system can be right for the wrong reasons in a way the score
makes visible --- a correct multiple-choice answer with a wrong paper, or a
correct paper with a locator pointing at the wrong page.

Our participation had an unusual shape and we want to be direct about it up
front. We built a conventional pipeline (dense retrieval over the released
metadata pool, cross-encoder reranking, an LLM selector, and a vision-capable
reader for tables and figures), it reached an official score of $0.4563$, and
then it stopped improving. Every subsequent gain --- to $0.7634$ --- came from
work that is better described as \emph{instrument reading} than as modelling:
studying the scorer, recovering the annotation conventions of the dataset from
the development split, and using controlled single-variable submissions to
identify which of our predictions were wrong.

We think the resulting findings are more useful to future participants than our
architecture is, so this paper is organised around them. Section~\ref{sec:metric}
decomposes the metric. Section~\ref{sec:system} describes the pipeline.
Section~\ref{sec:conventions} reports the annotation conventions we recovered.
Section~\ref{sec:attribution} describes score-guided attribution and its
limitations. Sections~\ref{sec:results}--\ref{sec:negative} give official
results, error analysis, and the heuristics we measured and discarded.

\section{The Metric Is Asymmetric}
\label{sec:metric}

The released evaluator computes
\begin{equation*}
\text{answer} = \frac{\text{mc} + \text{row\_f1} + \text{cell\_acc}}{3}
\end{equation*}
\begin{equation*}
\text{overall} = \frac{\text{paper\_f1} + \text{evidence\_f1} + \text{answer}}{3}
\end{equation*}
so paper F1 and evidence F1 each carry weight $1/3$, while multiple-choice
accuracy, table row F1 and table cell accuracy carry $1/9$ each. Two-thirds of
the score is therefore \emph{provenance}, not answering. This is the single most
consequential fact about the task and it is easy to miss: an intuition that
answer accuracy is the objective leads to effort in the wrong place. On our
first submission, multiple-choice accuracy was already $0.68$ while evidence F1
was $0.3587$.

The evidence metric deserves a second look. Predictions and gold are compared as
sets of coarse keys
\[
(\textit{paper\_id},\ \textit{source\_type},\ \textit{page},\ \textit{object\_id})
\]
where \textit{object\_id} is the normalised visible identifier for
\texttt{table}, \texttt{figure}, \texttt{equation\_algorithm} and
\texttt{citation\_context}, and is \emph{empty} for \texttt{text\_span}. Three
consequences shaped our work. (i) A \texttt{text\_span} item matches on paper and
page alone, so it is the cheapest correct guess when the type is uncertain.
(ii) Conversely, citing the right page with the wrong \textit{source\_type}
scores zero --- ``we found the answer'' is not the same as ``we located it as the
annotator did''. (iii) Because scoring is set-based and macro-averaged per
question, emitting $N$ predictions against $G$ gold items with $C$ correct gives
$F_1 = 2C/(G+N)$, which makes the cost of hedging exactly computable rather than
a matter of taste. We use this repeatedly.

The table metric has a similar structure. Row F1 is a set F1 over tuples of the
row-key columns, so duplicate row keys silently collapse in the evaluator's
dictionary (last one wins). Cell accuracy iterates over \emph{gold} rows only:
cells belonging to a predicted row whose key does not match gold are never
looked at, and unmatched gold rows count their cells as wrong. Cell accuracy is
therefore bounded by row \emph{recall} and is completely indifferent to row
precision --- which means an extra candidate row can cost row F1 but can never
cost cell accuracy. This asymmetry is the basis for several of our decisions in
Section~\ref{sec:conventions}.

\section{System}
\label{sec:system}

\paragraph{Retrieval and selection.}
The released pool contains 27{,}487 papers with title, abstract, venue, and a
PDF URL. We embed question and title--abstract text, retrieve a candidate list,
rerank with a cross-encoder, and pass the top candidates as a numbered list of
titles and abstracts to an LLM selector that returns indices. Enabling
reranking moved paper F1 from $0.410$ to $0.490$ on a development family, and the
LLM selector plus a cap of three papers per question was our best configuration
(\texttt{-{}-llm-select -{}-visual-table -{}-max-papers 3}).

We measured shortlist recall separately to know whether selection or retrieval
was the bottleneck: on a test-like development family, recall of the gold paper
was $0.846$ at rank 1 and $0.962$ at rank 20, and \emph{flat} out to rank 200.
Widening the shortlist was therefore pointless, and we did not pursue it.

\paragraph{Reading.}
Selected PDFs are fetched and read locally with PyMuPDF for page text, caption
indices, and page rasterisation. Answer generation is a per-question prompt
carrying the question, the requested answer types, the table schema when one is
supplied, and page text or rendered page images. Table questions use a two-stage
scheme in which row keys are extracted from the question alone before values are
read, guarded by a check that every extracted key actually occurs in the question
--- without that guard the stage fabricated entities and cost $0.129$ overall.

\paragraph{Audit loop.}
The component that produced most of our score is not a model but a procedure.
Each candidate submission is checked question by question against the source
PDFs: does the cited page contain the claim, does the cited object id exist in
that page's captions, does every emitted paper mention an entity the question
names, does the chosen multiple-choice option match the row and column the
question actually asks about. We automated the checks that could be automated
(Section~\ref{sec:results}) and performed the adjudication with an LLM agent
reading the rendered pages. This is a human-in-the-loop system in the sense that
matters for interpreting our score, and we return to it in
Section~\ref{sec:limitations}.

\section{Recovered Annotation Conventions}
\label{sec:conventions}

The development split has 55 examples with gold papers, evidence and answers.
Used as a validation set for tuning it is too small to separate configurations.
Used as a corpus describing \emph{how the annotators behaved}, it was the most
valuable resource in the task. Every statement below is a count over that split;
Table~\ref{tab:conventions} collects them with their denominators, because
several rest on small samples and we would rather a reader see that than take
them on trust.

\begin{table}[t]
\centering\small
\setlength{\tabcolsep}{4pt}
\begin{tabular}{p{5.1cm}r}
\toprule
Convention (counted on the 55 dev questions) & support \\
\midrule
Exactly one evidence key per gold paper            & 43/55 \\
\;\;of which $(1,1)$ / $(4,4)$ patterns            & 24 / 19 \\
Gold table rows $=$ gold paper count (multi-paper) & \textbf{8/8} \\
Question names a figure $\Rightarrow$ gold has \texttt{figure} & \textbf{10/10} \\
Question names a table $\Rightarrow$ gold has \texttt{table}   & 3/3 \\
Question asks an equation/loss $\Rightarrow$ gold has \texttt{equation\_algorithm} & 4/7 \\
Gold pairs a \texttt{text\_span} with a same-page object & 1/64 \\
Gold evidence page is the \emph{earliest} occurrence of the value & 10/38 \\
\;\;middle / latest                                & 22 / 6 \\
\bottomrule
\end{tabular}
\caption{Annotation conventions recovered from the development split, with the
counts they rest on. The two strongest (rows 3--4) were the most productive on
test; the last two are why we did \emph{not} ship a positional page-rewrite rule
or a same-page \texttt{text\_span} pad.}
\label{tab:conventions}
\end{table}

\paragraph{One evidence key per paper.}
Comparing the number of distinct coarse evidence keys to the number of gold
papers gives $(1,1)$ 24 times, $(4,4)$ 19 times, plus $(3,3)$ and $(9,9)$; 43 of
55 questions have exactly one key per paper. Gold evidence is therefore
\emph{provenance per paper}, not an exhaustive list of supporting locations. We
sized our evidence sets accordingly.

\paragraph{Source type follows the question's wording.}
Conditioning gold's type on the question's vocabulary is strikingly predictive.
Of the questions whose wording names a figure, plot or chart, \textbf{10 of 10}
have gold containing \texttt{figure}; for wording that names a table, 3 of 3;
for wording that asks for an equation, expression, loss or objective, 4 of 7
contain \texttt{equation\_algorithm}, and 2 of those 4 also carry
\texttt{text\_span}. Acting on this was our single largest evidence gain: we had
been emitting \emph{no} \texttt{equation\_algorithm} items at all. Adding them
alongside \texttt{text\_span} on the six equation-form questions coincided with
official evidence F1 moving from $0.6610$ to $0.7121$; that submission also
re-aimed two evidence pages, so the whole delta is not attributable to the type
change alone. Where the type is genuinely mixed in gold,
carrying both is the right arithmetic: against a single gold item, two
predictions with one hit score $2/3$, where committing to one scores $0.4$ in
expectation at $P(\text{equation})=0.6$.

\paragraph{Gold table rows equal gold paper count.}
For every multi-paper table question in the development split --- $(4,4)$ six
times, $(3,3)$, $(9,9)$ --- gold has exactly one row per paper, 8 of 8. This is a
hard structural constraint, and it does more than cap row F1: if gold has two
rows, then gold's two keys cannot be two of our four \emph{quantity descriptors},
because a two-row table cannot key its rows by four separate quantities. Three of
our test predictions violated the constraint and were, by this argument, scoring
zero.

\paragraph{Row keys reproduce printed labels.}
Where the question names entities, gold keys are those names lowercased, keeping
qualifying parentheticals (\texttt{ecm-xl (102.4m)}). Where the question implies
variants of one paper, gold uses the labels \emph{printed in that paper's own
table} --- one development question keys its rows \texttt{w/ ground truth} and
\texttt{w/ cube rcnn}. Checking every test row key against its source paper this
way found four mismatches, of which the decisive one was a retrieval question
whose gold keys carry the paper's \texttt{+Fake2} suffix rather than the bare
dataset names we had used. Correcting it took that question's row F1 from $0$ to
$1.0$ and its cell accuracy from $0$ to $1.0$ in one step.

\paragraph{Evidence pages are not the first mention.}
Gold's page histogram peaks at pages 6--7 (65 of 149 items) with only 5 of 149 on
page 1, while our predictions put 40\% of their items on pages 1--3. The obvious
correction is wrong, though: among gold items whose answer value appears on more
than one page, gold is the earliest occurrence 10 times, somewhere in the middle
22 times, and the latest 6 times. There is no positional rule to learn, only a
semantic one. We used the statistic as a triage signal for manual review and
explicitly did not ship it as a rewrite (Section~\ref{sec:negative}).

\section{Score-Guided Attribution}
\label{sec:attribution}

The task permits five evaluator submissions per day. Because each macro metric is
a mean of a small number of rational per-question values, a submission that
changes one thing often has exactly one arithmetic explanation, and the feedback
identifies \emph{which} of our predictions were wrong rather than merely whether
the file got better.

Two examples. Table row F1 sat at exactly $0.340476$ for five consecutive
submissions, so when two questions' row keys were changed and the macro moved
by $+0.011111$ --- that is, $+0.2333$ over the 21-question sum --- the shift
decomposed uniquely as one question going from zero matched rows to one while its
row count went from 2 to 4 ($2\cdot1/6-0=+0.3333$) and a second question's added
candidate missing ($2\cdot1/5-2\cdot1/4=-0.1$). That told us both that gold's
\texttt{paper} row keys are short names rather than full metadata titles, and
that one specific candidate string was wrong. Second,
\texttt{table\_cell\_accuracy\_micro} is $\textit{cell\_correct}/\textit{cell\_total}$
and \textit{cell\_total} is a property of gold: reading it as an exact fraction
gives $23/87$ and later $31/87$, so gold has exactly 87 comparable cells and the
$+8$ was one question's eight cells.

The same instrument disciplined us. A later submission gained $+0.3333$ where we
had predicted $+0.5833$; because three of the changed questions had identical key
sets before and after, their contributions had to cancel, which forced the
conclusion that a change we had reasoned was ``free'' had in fact discarded a
matching key. We reverted it.

We flag clearly that this is adaptation to the held-out split. It is within the
stated submission rules, and we report it because it materially explains our
final number; but scores obtained this way should not be read as an estimate of
how the underlying system would perform on unseen questions.

\section{Official Results}
\label{sec:results}

Table~\ref{tab:results} reports official evaluator output for our scored
submissions on \texttt{littraceqa-test} (71 questions). All figures are as
returned by the evaluator; no numbers in this paper are computed by us from
private labels.

\begin{table*}[t]
\centering\small
\begin{tabular}{lrrrrrr}
\toprule
Submission & paper F1 & evid.\ F1 & MC & row F1 & cell acc & \textbf{overall} \\
\midrule
initial pipeline                     & 0.6324 & 0.3587 & 0.68 & 0.3221 & 0.1310 & 0.4563 \\
\;+ LLM selector, visual tables      & 0.7991 & 0.4737 & 0.78 & 0.2738 & 0.0952 & 0.5519 \\
\;+ hand-authored table row keys     & 0.7991 & 0.4737 & 0.78 & 0.3405 & 0.0952 & 0.5602 \\
\;+ answer audit vs.\ PDFs           & 0.8554 & 0.5347 & 0.84 & 0.3405 & 0.1984 & 0.6166 \\
\;+ audit, continued                 & 0.8789 & 0.5606 & 0.86 & 0.3405 & 0.2103 & 0.6366 \\
\;+ 7 paper-set corrections          & 0.9704 & 0.6610 & 0.94 & 0.3405 & 0.2103 & 0.7095 \\
\;+ \texttt{equation\_algorithm} evidence & 0.9704 & 0.7121 & 0.96 & 0.3405 & 0.2103 & 0.7287 \\
\;+ figure-counting answers          & 0.9704 & 0.7121 & 0.98 & 0.3516 & 0.2103 & 0.7322 \\
\;+ short-name table row keys        & 0.9704 & 0.7121 & 0.98 & 0.4283 & 0.2262 & 0.7425 \\
\;$-$ rows proven not to match       & 0.9704 & 0.7121 & 0.98 & 0.4675 & 0.2262 & 0.7468 \\
\;+ batched cell rewrites            & 0.9704 & 0.7192 & 0.98 & 0.4675 & 0.2500 & 0.7518 \\
\;+ printed-label row keys           & 0.9704 & 0.7262 & 0.98 & 0.4873 & 0.2976 & 0.7617 \\
\textbf{final submission}            & \textbf{0.9704} & \textbf{0.7262} & \textbf{0.98} & \textbf{0.5032} & \textbf{0.2976} & \textbf{0.7634} \\
\bottomrule
\end{tabular}
\caption{Official evaluator results on the LitTraceQA held-out test split.
Rows are in submission order; each line names the change that distinguishes it
from the line above. Paper F1 and evidence F1 carry weight $1/3$ each in
\emph{overall}; MC, row F1 and cell accuracy carry $1/9$ each. Components are the
evaluator's reported values rounded to four decimals, so recomputing
\emph{overall} from them can differ in the last digit. Row two is our best fully
automated submission; every later row involves per-question intervention.}
\label{tab:results}
\end{table*}

\paragraph{Where the score came from.}
Between our first and last submission, paper F1 rose by $0.338$ and evidence F1
by $0.368$; multiple-choice accuracy rose by $0.30$ but contributes at one third
the weight. Table row F1 and cell accuracy remain our weakest components at
$0.5032$ and $0.2976$, and together they account for the majority of our
remaining headroom.

\paragraph{Automated checks.}
Four checks ran over every candidate file and are the reusable part of our
tooling. (1) An \emph{object-id verifier} regexes captions out of each page and
confirms the cited \texttt{table\_id}/\texttt{figure\_id} exists on the cited
page; 48 of 49 of our ids passed, and the one failure was a real error (a
retention-score column we had attributed to Table~2 when its caption sits on the
next page). (2) A \emph{name-presence check} asks whether each emitted paper's
text ever mentions an entity the question names --- a paper that never writes
``ERASE'' is not the ERASE paper. (3) The same check against all 27{,}487 titles,
which finds papers that were never retrieved and so could not be checked any
other way. (4) A \emph{dead-run guard} that fails a submission containing
all-null table rows; this caught two files silently produced after our API quota
was exhausted mid-run, with no error in the trace.

\section{Error Analysis}
\label{sec:errors}

\paragraph{The dominant failure is the right paper and the wrong row.}
Of the answer errors we found and corrected, almost none are hallucinations.
They are values genuinely printed in the cited paper, one row or one column away
from the one the question asks about: the un-finetuned base-model row instead of
the fine-tuned one; \texttt{RPEt} where \texttt{RPEr} was asked; the \texttt{MC1}
column where the question says \texttt{MC3}; inter-annotator agreement where the
question says intra-annotator; a \emph{Pearson} column read as a Kendall one;
the \emph{None} baseline row's accuracy reported as the proposed method's
macro-F1. This is why three separate value-presence triages we built returned
mostly false positives --- every candidate value \emph{is} in the paper --- and
why reading row and column headers is the only check that works.

\paragraph{Wrong papers hide behind plausible titles.}
Seven of our paper sets were wrong and every one was found by the presence
checks rather than by re-running selection. Two required fetching a PDF that had
never been retrieved, so no text-based check could have reached them, and one was a question
answered ``all three papers'' while citing a paper that was none of the three the
options named.

\paragraph{Figures need rasterising.}
Four questions ask for counts that exist only in a figure. Rendering those pages
at $2\times$ and reading them settled all four and corrected one answer: a paper
states in prose that it classifies models into \emph{eight} categories, and its
pie chart labels agree, where we had answered six.

\paragraph{Grader-invented strings are not recoverable.}
Where a table's row-key column is a descriptor the annotator invented
(\texttt{quantity}, \texttt{setting}, \texttt{attribute}, \texttt{metric}), we
never matched it. One such cell was attempted in four different forms --- the
verbatim question clause, a bare noun phrase, the same with a ``number of''
prefix, and the source paper's own noun --- and all four missed. We report this as
a limit of the approach rather than as an unexplored lead.

\section{Negative Results}
\label{sec:negative}

We list heuristics that seemed reasonable, were measured, and were dropped. We
think these are the most transferable part of the paper.

\begin{itemize}\itemsep2pt
\item \textbf{Padding paper sets.} From one submission's precision and recall
alone, the gold set sizes on test are derivable as 1 or 2 (3 and 4 are
arithmetically impossible), and our selector is size-calibrated: an
uncalibrated selector cannot reach the observed precision even if its second
pick were always gold. Padding singletons was therefore refused.
\item \textbf{Full-text indexing.} Content-defined multi-paper questions look
like they need body text, but only 4\% of the relevant development family is
answerable from body text alone; the shortlist was already saturated
(Section~\ref{sec:system}).
\item \textbf{Re-aiming evidence at later pages.} Motivated by the page
histogram, measured at $+0.013$ evidence F1 on development --- below our
shipping threshold --- and explained by the middle-position statistic.
\item \textbf{Carrying both row-key spellings everywhere.} Measured at
$-0.020$ row F1 on development, because there the bare keys are usually already
right and the alternative never hit. We restricted hedging to keys we could
\emph{prove} wrong.
\item \textbf{Adding a \texttt{table} item wherever a text-span page has a table
caption.} Initially recorded as ``no effect''; on re-examination the experiment
added only 4 items and all 4 landed on questions already scoring zero, so it was
\emph{uninformative rather than refuted}. We did not ship it, and we correct the
earlier reading here.
\item \textbf{Retyping \texttt{text\_span} to \texttt{table}.} Of 44 wrong
evidence items examined, 41 were page-level errors and only 2 were type errors,
so a blanket retype had no target.
\end{itemize}

We also record two errors in our own instrumentation, because both produced
conclusions we briefly believed. A multiple-choice triage stripped whitespace
before matching numbers, which glues adjacent table cells together, so
\texttt{16,825,000 5,800,000} became one token and both halves were reported
absent; every ``value not in this paper'' verdict it produced was worthless until
whitespace was collapsed rather than removed. And a prediction-file diff compared
a key that prediction files do not carry, so both sides read \texttt{None} and
the diff reported ``0 of 71 identical'' when the true answer was 26.

\section{Setup and Disclosed Resources}
\label{sec:setup}

\paragraph{Data.} Released LitTraceQA metadata pool (27{,}487 papers),
development split (55), and input-only test split (71), from the official
Hugging Face dataset (CC BY-NC 4.0). PDFs were fetched from the publisher URLs
in the released metadata. We used no additional annotated data.

\paragraph{Models and APIs.} Hosted \texttt{gemini-flash-lite-latest} and
sibling Gemini models on the free tier for selection and reading; a locally
served Qwen3 through an OpenAI-compatible \texttt{llama.cpp} endpoint for
text-only selection experiments, added specifically because the hosted free-tier
quota (500 requests/day) is too small to score one development split twice; a
cross-encoder reranker and a sentence embedding model for retrieval. The audit
loop described in Section~\ref{sec:system} was carried out by an LLM agent
(Anthropic Claude) reading page text and rasterised pages, with all adjudications
recorded against a specific page of a specific PDF.

\paragraph{Software and compute.} Python, PyMuPDF for PDF text/caption/raster
extraction, the organisers' released \texttt{evaluate.py} and
\texttt{validate\_submission.py} for local checks. All work ran on a single
workstation with no GPU training; we fine-tuned nothing.

\paragraph{Reproducibility.} We release the pipeline, the four automated
verifiers, the experiment scripts, and the reports recording each measurement
including the negative ones. Prediction files and logs are not committed; the
reports are the record.

\section{Limitations}
\label{sec:limitations}

The most important limitation is the nature of our final submission. Our best
\emph{fully automated} run --- pipeline output with no per-question intervention
--- scores $\mathbf{0.5519}$ (row two of Table~\ref{tab:results}); the reported
$0.7634$ additionally includes per-question auditing against source PDFs and
corrections attributed using leaderboard feedback. The audit is not a model that
would run on new questions, and the attribution is adaptation to this particular
held-out split. A reader interested in what an automated system achieves on
LitTraceQA should read the $0.5519$ row, not the last one. We report both
prominently for exactly this reason.

Second, our conventions in Section~\ref{sec:conventions} are counts over 55
development questions, and several rest on small denominators (3 of 3 for the
table conditional, 8 of 8 for the row-count rule). They were predictive on test,
which is evidence, but they are not established at any conventional level of
confidence.

Third, our weakest components are the table ones, and our analysis says why we
could not fix them: matching a row key that an annotator invented is not a
reasoning problem with a findable answer. This is arguably a property of the
task's table format rather than of our system, and it interacts badly with the
metric --- because cell accuracy is gated on row-key recall, a question whose key
wording we cannot reproduce contributes zero to two of the five reported metrics
no matter how correct the extracted values are.

Fourth, we did not submit to the optional \texttt{littraceqa-test-extra}
diagnostic track, so we cannot report how these findings transfer to its larger
question pool.

\section{Ethics}

Papers were retrieved from publisher-provided URLs in the released metadata and
used only for evidence extraction. The dataset is CC BY-NC 4.0 and we used it
within those terms. Our system's purpose is to make claims about scientific
literature checkable by pointing at a specific page of a specific paper, which we
consider a beneficial direction; the corresponding risk is that a confidently
formatted citation is more persuasive than an uncited claim even when wrong, and
our own error analysis shows the dominant failure mode is a citation that is
almost right --- correct paper, adjacent row. Systems of this kind should surface
the located span for human confirmation rather than presenting a grounded answer
as settled.

\section{Conclusion}

Our LitTraceQA submission reached an official $0.7634$, and the useful part of
the story is where the gain came from: two-thirds of the metric is provenance
rather than answering, gold's annotation conventions are recoverable from 55
examples, and a macro metric over 71 questions is a precise enough instrument
that a controlled submission can prove an individual prediction wrong. We also
report the heuristics this process refuted, two bugs in our own verifiers, and
the extent to which our final number reflects auditing rather than a system.
For future participants we would summarise the task in one sentence: locate the
evidence the way the annotator located it, and check the row and column headers
before believing a value you have found.

\section*{Acknowledgements}

We thank the GroundLM 2026 organisers for the released evaluator and validator,
which made the analysis in this paper possible.

\begin{thebibliography}{9}
\bibitem[LitTraceQA(2026)]{littraceqa}
LitTraceQA dataset (CC BY-NC 4.0).\\
\url{https://huggingface.co/datasets/LitTraceQA/LitTraceQA}

\bibitem[GroundLM(2026)]{groundlm}
GroundLM 2026 Workshop, EMNLP 2026 --- shared task description, evaluator and
submission rules.\\
\url{https://groundlm.github.io/grouplm_emnlp2026/shared-tasks.html}

\bibitem[GroundLM Evaluator(2026)]{evaluator}
GroundLM 2026 online evaluator (source of record for all held-out scores
reported here).\\
\url{https://huggingface.co/spaces/YimuWang/GroundLM-2026-Eval}

\bibitem[PyMuPDF(n.d.)]{pymupdf}
PyMuPDF --- Python bindings for MuPDF, used for page text, caption indices and
page rasterisation.\\
\url{https://pymupdf.readthedocs.io}

\bibitem[ACL(2026)]{aclstyle}
ACL style files and ACL PubCheck.\\
\url{https://github.com/acl-org/acl-style-files}
\end{thebibliography}

\end{document}
